\section{Introduction}
\label{sec:intro}

% P1: task and why token cost matters
Language models read tokens, not characters, and the cost of a request, its latency and the
amount of text that fits in the context window are all counted in tokens. Byte-pair
encoding (BPE) \citep{gage1994,sennrich2016} builds its vocabulary by repeatedly merging the
most frequent pair of adjacent symbols in a training corpus, so the merges follow the
languages that dominate the corpus, and the same content needs many more tokens in some
languages than in others, up to 15 times more for deployed tokenizers \citep{petrov2023}.
In our setting, a Hindi request to a model with an English-heavy tokenizer costs 21\% more
tokens than its English translation (35.0 vs.\ 29.0 per sentence).

% P2: existing methods sit at two ends, the space between is unprincipled
Existing tokenizers answer this disparity at two opposite ends. Classical BPE minimises the
total number of tokens and accepts whatever disparity results; Parity-aware BPE
\citep{foroutan2026} gives every merge to the language that is currently worst off, which
equalises the cost across languages at some loss of total compression. To reach points in
between, \citet{foroutan2026} propose a hybrid schedule that switches from the classical
rule to the parity rule after a fixed share of merges, and \citet{kanjirangat2026} extend
the parity rule to the $N$ worst-off languages. Both are heuristic interpolations: neither
is derived from an objective that describes the trade-off, so a practitioner cannot state
what the intermediate tokenizers optimise or choose a point on the trade-off deliberately.

% P3: our proposal
We treat vocabulary learning as a multi-objective problem over the vector of per-language
token costs, scalarise it with a weighted power mean of exponent $p$, and show that the
greedy step of this objective is a pair count in which each language is weighted by its
current cost raised to the power $p-1$ (Section~\ref{sec:method}). At $p=1$ the rule is
exactly classical BPE, and as $p\to\infty$ it becomes the rule of Parity-aware BPE; we call
the tokenizers in between \emph{$p$-fair BPE}. We further prove that the score of an
existing pair never increases during training, which allows exact training with a lazy
priority queue \citep{minoux1978} at about 20 times the speed of full rescoring.

% P4: what we find
We evaluate $p$-fair BPE on Hindi, Marathi, Gujarati and English with five samples of
training text and a paired bootstrap on FLORES-200 \citep{nllb2022}, and on 14 languages in
ten scripts at natural data balance with three. Raising $p$ moves the tokenizer monotonically from the classical end to the parity end in
every sample, and at $p=8$ the Gini coefficient of per-language cost falls by 40\% for a
0.6\% loss in compression. The front coincides with that of the hybrid schedule in both settings, where
the two differ by at most 0.003 in Gini at equal compression, so $p$-fair BPE matches the
hybrid and does not dominate it; only at a 128k vocabulary is it fairer than every hybrid
point, by 0.002 to 0.006 in five of five samples. Nevertheless, the two methods distribute
the cost differently: for every $p$ we test, no language that costs more than the classical
mean becomes more expensive, whereas the hybrid raises Marathi with four languages (33.26
vs.\ 32.47 tokens per sentence at 75\% classical merges) and Malayalam and Gujarati with
fourteen. Three further factors decide the outcome more than the merge rule does: a per-byte cost
inverts the parity rule across scripts, sending 99\% of its merges to English; the Gini
coefficient of classical BPE ranges from 0.019 with equal data to 0.051 with eight times as
much English; and intermediate $p$ loses equity as the vocabulary grows from 32k to 256k,
whereas the parity rule does not.

\noindent\textbf{Contributions.} (i) We formulate fair vocabulary learning as the
minimisation of a power mean of per-language token costs and show that its greedy step
contains classical BPE ($p=1$) and Parity-aware BPE ($p\to\infty$) as limits of one rule.
(ii) We prove that pair scores are non-increasing under this rule, which gives an exact
lazy-evaluation trainer about 20 times faster than full rescoring. (iii) We trace the
compression--parity front with error bars on four languages and on 14 languages in ten
scripts, and show that $p$-fair BPE matches the hybrid front in both settings while never
raising the cost of an under-served language above its classical cost. (iv) We show that
the cost unit, the data balance and the vocabulary size decide the size and even the
direction of the disparity, that Worst-$N$ parity and the moving window do not
interpolate, and that a convex vocabulary allocation across scripts is a weaker lever than
the merge rule.